\section{Манга и раннее аниме: как выразительная условность победила реализм}

Прямой мост между театральной маской и современным аниме-архетипом — Осаму Тэдзука. Его "Тэцуван Атому" ("Astro Boy"), манга 1952 года и аниме-сериал 1963-го, задал визуальный код, который до сих пор работает: непропорционально большие глаза, упрощённые черты лица, минимум линий там, где реализм потребовал бы десятки. Тэдзука сам не скрывал источник: он открыто восхищался диснеевской "Бэмби" и мультфильмами студии Fleischer и прямо заимствовал оттуда крупный, эмоционально нагруженный глаз.

Но заимствование вкуса — не единственная причина, по которой код прижился именно так и именно тогда. Решающим фактором была экономика, а не только эстетика. Телевизионное аниме 1960-х делалось на бюджетах, несопоставимых с диснеевской полнометражной анимацией: не 24 кадра в секунду, а восемь, часто меньше, с придержанными статичными кадрами и переиспользуемыми фазами движения. В таких условиях подробное реалистичное лицо, которое требует плавной мимики через десятки промежуточных фаз, попросту не работает — оно либо дорого, либо на статичном кадре выглядит мёртвым.

Вот где решение снова, как и в 1.1, оказывается структурно тем же самым, только найденным заново под новым ограничением: если нельзя непрерывно менять всё лицо, нужно сосредоточить всю эмоциональную нагрузку на минимальном участке, который можно перерисовать одним движением руки. Глаз идеально подходит на эту роль — большая площадь, максимальная считываемость формы и направления взгляда, и всего одна деталь, которую нужно изменить, чтобы полностью сменить эмоциональный регистр всего лица: сузить — угроза, расширить — потрясение, притенить — скрытность. Рот и брови можно оставить почти неизменными; всю работу выполняет глаз. Это тот же принцип, который МакКлауд опишет как amplification through simplification (0.3), только здесь он продиктован не эстетической теорией, а прямым производственным расчётом: где выгоднее всего разместить всю эмоциональную нагрузку персонажа, если перерисовывать можно немного и редко.

К концу 1960-х этот код уже не воспринимался как техническое ограничение — он стал языком, а глаз — первым в истории аниме модульным элементом эмоции: деталью, которую можно менять независимо от остального лица и которая одна несёт полное сообщение. Это ранняя, ещё безымянная версия того, что в 1.5 получит имя moe-элемента, а индустрия превратит в отдельный рынок распознаваемых, комбинируемых деталей. Большие глаза читались как открытость и искренность независимо от сюжета; их сужение или закрытие тенью стало отдельным, самостоятельным сигналом драмы — задолго до того, как кто-то начал сознательно собирать персонажей из подобных деталей.
